\documentclass[12pt,letterpaper,twoside,openright]{book}

% --- Packages ---
\usepackage[utf8]{inputenc}
\usepackage[T1]{fontenc}
\usepackage{amsmath,amssymb,amsfonts,amsthm}
\usepackage{mathtools}
\usepackage{geometry}
\usepackage{hyperref}
\usepackage{cleveref}
\usepackage{enumitem}
\usepackage{booktabs}
\usepackage{microtype}
\usepackage{natbib}
\usepackage{tikz-cd}
\usepackage{tikz}
\usepackage{graphicx}
\usepackage{float}
\usepackage{xcolor}
\usepackage{fancyhdr}
\usepackage{longtable}
\usepackage{array}
\usepackage{makeidx}
\usepackage{tcolorbox}
\usepackage{listings}
\usepackage{algorithm}
\usepackage{algpseudocode}
\usepackage{epigraph}
\usepackage{framed}

\geometry{margin=1in, inner=1.25in, outer=1in}

\hypersetup{
    colorlinks=true,
    linkcolor=blue!70!black,
    citecolor=green!50!black,
    urlcolor=blue!60!black
}

% --- Index ---
\makeindex

% --- Header/Footer ---
\pagestyle{fancy}
\fancyhf{}
\fancyhead[LE]{\small\textit{Philosophy Engineering}}
\fancyhead[RO]{\small\textit{\nouppercase{\rightmark}}}
\fancyfoot[C]{\thepage}
\renewcommand{\headrulewidth}{0.4pt}

% --- Theorem environments ---
\theoremstyle{plain}
\newtheorem{theorem}{Theorem}[chapter]
\newtheorem{lemma}[theorem]{Lemma}
\newtheorem{corollary}[theorem]{Corollary}
\newtheorem{proposition}[theorem]{Proposition}
\theoremstyle{definition}
\newtheorem{definition}[theorem]{Definition}
\newtheorem{example}[theorem]{Example}
\newtheorem{axiom}{Axiom}[chapter]
\newtheorem{principle}[theorem]{Principle}
\theoremstyle{remark}
\newtheorem{remark}[theorem]{Remark}

% --- Pedagogical environments ---
\tcbuselibrary{breakable,skins}

\newtcolorbox{keyinsight}[1][]{
    colback=blue!5!white, colframe=blue!60!black,
    fonttitle=\bfseries, title=Key Insight, #1
}

\newtcolorbox{practitioner}[1][]{
    colback=green!5!white, colframe=green!50!black,
    fonttitle=\bfseries, title=Practitioner Note, #1
}

\newtcolorbox{historicalnote}[1][]{
    colback=orange!5!white, colframe=orange!50!black,
    fonttitle=\bfseries, title=Historical Note, #1
}

\newtcolorbox{pitfall}[1][]{
    colback=red!5!white, colframe=red!60!black,
    fonttitle=\bfseries, title=Common Pitfall, #1
}

\newtcolorbox{casestudy}[1][]{
    colback=yellow!5!white, colframe=yellow!50!black,
    fonttitle=\bfseries, title=Case Study, breakable, #1
}

% --- Code listing style ---
\lstset{
    basicstyle=\ttfamily\small,
    keywordstyle=\color{blue!70!black}\bfseries,
    commentstyle=\color{green!50!black}\itshape,
    stringstyle=\color{red!60!black},
    frame=single,
    breaklines=true,
    numbers=left,
    numberstyle=\tiny\color{gray},
    captionpos=b
}

% --- Custom commands ---
\newcommand{\R}{\mathbb{R}}
\newcommand{\N}{\mathbb{N}}
\newcommand{\calD}{\mathcal{D}}
\newcommand{\calT}{\mathcal{T}}
\newcommand{\calU}{\mathcal{U}}
\newcommand{\calB}{\mathcal{B}}
\newcommand{\calL}{\mathcal{L}}
\newcommand{\calM}{\mathcal{M}}
\newcommand{\calH}{\mathcal{H}}
\newcommand{\Lag}{\mathcal{L}}
\newcommand{\BF}{\mathrm{BF}}

% --- Exercise counter ---
\newcounter{exercise}[chapter]
\renewcommand{\theexercise}{\thechapter.\arabic{exercise}}
\newenvironment{exercise}[1][]{%
    \refstepcounter{exercise}%
    \par\medskip\noindent\textbf{Exercise~\theexercise\if\relax\detokenize{#1}\relax\else\ (#1)\fi.}\itshape\enspace
}{\par\medskip}

% ============================================================
\begin{document}

% ============================================================
% FRONT MATTER
% ============================================================
\frontmatter

% --- Half title ---
\begin{titlepage}
\vspace*{3cm}
\begin{center}
{\Huge\bfseries Philosophy Engineering}
\end{center}
\vfill
\end{titlepage}

\cleardoublepage

% --- Full title page ---
\begin{titlepage}
\centering
\vspace*{2cm}

{\Huge\bfseries Philosophy Engineering\par}
\vspace{0.8cm}
{\Large\itshape Translating Philosophical Commitments into\\Computationally Enforceable Constraints\par}

\vspace{2cm}

{\large Andrew H. Bond\par}
\vspace{0.3cm}
{\normalsize Department of Computer Engineering\\
San Jos\'{e} State University\par}

\vspace{2cm}

{\normalsize First Edition\par}
\vspace{0.3cm}
{\normalsize 2026\par}

\vfill

% Publisher placeholder
{\small\textit{[Publisher]\par}}

\end{titlepage}

\cleardoublepage

% --- Copyright page ---
\thispagestyle{empty}
\vspace*{\fill}
\noindent Copyright \copyright\ 2026 Andrew H. Bond. All rights reserved.

\vspace{1em}
\noindent ISBN: [TBD]

\vspace{1em}
\noindent\textbf{AI-Assisted Writing Disclosure:} Portions of this manuscript were drafted and refined with AI assistance (Claude, Anthropic). All technical content, mathematical results, and claims were conceived, verified, and approved by the author.

\vspace{1em}
\noindent Typeset in \LaTeX.

\vspace{1em}
\noindent Library of Congress Cataloging-in-Publication Data: [TBD]
\cleardoublepage

% --- Epigraph ---
\thispagestyle{empty}
\vspace*{3cm}
\begin{center}
\textit{``Differences that do not matter should not make a difference.''}

\vspace{0.5cm}
--- The Epistemic Invariance Principle
\end{center}
\vfill
\cleardoublepage

% --- Preface ---
\chapter{Preface}

\section*{Who This Book Is For}

This textbook is written for three audiences:

\begin{enumerate}
    \item \textbf{Graduate students} in computer science, AI/ML, philosophy, and interdisciplinary programs who want a rigorous introduction to the formal foundations of value-aligned system design.
    \item \textbf{Practicing engineers} building safety-critical AI systems who need actionable methods for ensuring that their systems respect epistemic and ethical norms---not as aspirational guidelines, but as testable, enforceable contracts.
    \item \textbf{Researchers} in AI alignment, formal ethics, and mathematical philosophy who seek a unified mathematical treatment connecting these fields.
\end{enumerate}

\section*{Prerequisites}

The book assumes:
\begin{itemize}
    \item Undergraduate linear algebra and multivariable calculus
    \item Basic set theory and proof techniques
    \item Introductory programming (Python examples use ErisML)
    \item Exposure to at least one of: philosophy of science, formal logic, or machine learning
\end{itemize}

\noindent No prior knowledge of differential geometry, gauge theory, or formal ethics is required. All necessary mathematics is developed from first principles.

\section*{How to Use This Book}

The book is organized in five parts. Parts~I--II are foundational and should be read in order. Parts~III--V can be read in any order depending on the reader's interests:

\begin{itemize}
    \item \textbf{One-semester graduate course}: Chapters 1--5, 8--10, 14--15, 19
    \item \textbf{AI alignment track}: Chapters 1--3, 8--12, 14--16
    \item \textbf{Mathematics track}: Chapters 4--7, 13, 17--18
    \item \textbf{Practitioner fast path}: Chapters 1--2, 10, 14--16, 19--20
\end{itemize}

Each chapter includes worked examples, exercises at three difficulty levels (\textbf{A}~=~routine, \textbf{B}~=~intermediate, \textbf{C}~=~research-level), and practitioner notes connecting theory to implementation.

\section*{Companion Software}

All code examples use ErisML v3.0, available at:
\begin{center}
\url{https://github.com/ahb-sjsu/erisml-lib}
\end{center}
Installation instructions and Jupyter notebooks for every chapter are provided in the companion repository.

\section*{Acknowledgments}

[To be written.]

\cleardoublepage

% --- Table of Contents ---
\tableofcontents
\cleardoublepage

% --- List of Figures, Tables, Algorithms ---
\listoffigures
\cleardoublepage
\listoftables
\cleardoublepage

% --- Notation ---
\chapter{Notation and Conventions}

\begin{longtable}{lp{10cm}}
\toprule
\textbf{Symbol} & \textbf{Meaning} \\
\midrule
$\calM$ & Moral manifold \\
$g_{ij}$ & Metric tensor on $\calM$ \\
$\calH$ & Harm functional \\
$\calL$ & Lagrangian density \\
$\Lag_\eta$ & Normative Lagrangian \\
$G$ & Gauge group (typically $D_4$) \\
$\pi: P \to \calM$ & Principal bundle \\
$A_\mu$ & Gauge connection \\
$F_{\mu\nu}$ & Field strength / curvature \\
$J^\mu$ & Conserved normative current \\
$\calT$ & Transformation set \\
$\equivX$ & Equivalence relation on input space $X$ \\
$\calB$ & Bond structure (objective relational data) \\
$\phi$ & Evaluative lens \\
$\beta(f, \calT)$ & Bond Index \\
$\Sigma_i$ & Stratum of a Whitney stratified space \\
$\R^9$ & Nine-dimensional moral evaluation space \\
\bottomrule
\end{longtable}

% ============================================================
% MAIN MATTER
% ============================================================
\mainmatter

% ============================================================
% PART I: FOUNDATIONS --- WHY PHILOSOPHY ENGINEERING?
% ============================================================
\part{Foundations: The Case for Philosophy Engineering}
\label{part:foundations}

% Chapter 1 ------------------------------------------------
\chapter{The Representation Dependence Crisis}
\label{ch:crisis}

\begin{epigraph}
\textit{``A distinction without a difference is no distinction at all.''}
\end{epigraph}

% Learning objectives
\begin{framed}
\noindent\textbf{Learning Objectives.} After completing this chapter, you will be able to:
\begin{enumerate}
    \item Identify representation dependence failures in deployed AI systems
    \item Distinguish representation dependence from bias and fairness concerns
    \item Explain why representation dependence is an epistemic (not merely statistical) failure
    \item State the core commitment of Philosophy Engineering in one sentence
\end{enumerate}
\end{framed}

\section{When Surface Form Changes the Verdict}
% AI systems that change output under paraphrase, reordering, formatting, etc.
% Concrete examples: medical diagnosis, legal assessment, content moderation
% Key point: these are not "bugs" --- they are symptomatic of a missing formal structure

\section{The Pragmatist Diagnosis}
% Peirce, James, Dewey: meaning is exhausted by practical consequences
% If two descriptions yield identical consequences in every scenario, they are equivalent
% A system sensitive to irrelevant representational choices is epistemically defective

\section{Beyond Bias: Representation vs.\ Content Failures}
% Bias concerns: the *content* of judgments (e.g., racial disparities)
% Representation dependence: the *form* of judgment itself
% A system can be unbiased and still epistemically broken
% These are orthogonal failure modes requiring different tools

\section{The Scalar Alignment Crisis}
% RLHF, constitutional AI, scalar rewards --- all assume moral evaluation is a scalar
% Information Monotonicity Theorem: scalar contraction is irreversibly lossy
% Specification gaming, reward hacking, value collapse as *mathematical* consequences
% Preview: Geometric Ethics provides the correct non-scalar structure

\section{What Philosophy Engineering Proposes}
% The four-step method: declare equivalences, declare invariances, enforce, audit
% Contracts are falsifiable, localizable, non-trivial
% Preview of the rest of the book

\begin{exercise}[A]
Find three examples of representation dependence in a publicly available language model. For each, identify the irrelevant surface change and the resulting output change.
\end{exercise}

\begin{exercise}[B]
A content moderation system flags ``I'm going to kill it at the presentation'' as violent but not ``I will perform excellently at the presentation.'' Is this a bias failure, a representation dependence failure, or both? Justify your answer using the distinctions from Section~1.3.
\end{exercise}

\begin{exercise}[C]
Formalize the claim ``representation dependence and bias are orthogonal failure modes.'' Construct a $2 \times 2$ matrix of systems that are (biased, representation-dependent), (biased, representation-independent), (unbiased, representation-dependent), and (unbiased, representation-independent). Provide concrete examples for each cell.
\end{exercise}


% Chapter 2 ------------------------------------------------
\chapter{What Is Philosophy Engineering?}
\label{ch:what-is-pe}

\begin{framed}
\noindent\textbf{Learning Objectives.} After completing this chapter, you will be able to:
\begin{enumerate}
    \item Define Philosophy Engineering and distinguish it from adjacent disciplines
    \item Describe the three-tier architecture (Substrate, Computation, Values)
    \item Identify the 16 knowledge areas and their tier assignments
    \item Explain the relationship to SWEBOK and why a new discipline is needed
\end{enumerate}
\end{framed}

\section{Definition and Scope}
% Philosophy Engineering = systematic translation of philosophical commitments into
% computationally enforceable constraints on autonomous systems
% Relationship to software engineering, AI/ML engineering, philosophy, AI ethics
% The disciplinary gap each leaves unfilled

\section{The Three-Tier Architecture}
% Tier I --- Substrate: math and physics providing representational vocabulary
% Tier II --- Computation: algorithms translating formal structures into code
% Tier III --- Values: normative disciplines defining objectives and constraints
% Diagram: the stack, with arrows showing dependencies

\section{The Sixteen Knowledge Areas}
% Table: KA number, name, tier, one-line description
% Substrate: Differential Geometry, Topology, Tensor Algebra, Group Theory,
%   Analytical Mechanics, Information Theory
% Computation: Algorithm Design, Formal Semantics, Formal Verification,
%   Network Topology, AI Alignment
% Values: Metaethics, Normative Ethics, Epistemology, Philosophy of Mind,
%   Political/Legal Philosophy

\section{Relationship to SWEBOK}
% IEEE/ACM SWEBOK as structural model
% How PE-BoK parallels and extends SWEBOK
% Competency levels: awareness, working knowledge, mastery

\section{Epistemic Stance}
% Pragmatism: mathematical structures as tools, not metaphysical commitments
% Epistemic status labels: Definition, Axiom, Theorem, Principle
% Falsifiability: all invariance claims are testable

\begin{exercise}[A]
For each of the 16 knowledge areas, identify one concrete way it contributes to the problem of building representation-independent AI systems.
\end{exercise}

\begin{exercise}[B]
Compare the SWEBOK knowledge area structure with the PE-BoK structure proposed here. Where do they overlap? Where do they diverge? What does the divergence tell us about the nature of Philosophy Engineering?
\end{exercise}


% Chapter 3 ------------------------------------------------
\chapter{A Brief History of Formal Ethics}
\label{ch:history}

\begin{framed}
\noindent\textbf{Learning Objectives.} After completing this chapter, you will be able to:
\begin{enumerate}
    \item Trace the historical arc from Aristotle through Kant to computational ethics
    \item Identify the key moments where formalization of ethics was attempted and why earlier attempts fell short
    \item Explain what changed with the availability of differential geometry and gauge theory
\end{enumerate}
\end{framed}

\section{Aristotle to Bentham: Systematizing Moral Reasoning}
% Aristotle's virtue theory as proto-formalization
% Bentham's felicific calculus: the first scalar reduction
% Mill's qualitative hedonism: the first hint that scalars are insufficient

\section{Kant's Formalism}
% The categorical imperative as an invariance test
% Universalizability = representation independence (avant la lettre)
% What Kant got right; what he lacked (no manifold, no metric)

\section{The Analytic Turn: Moore, Hare, and the Is-Ought Gap}
% Moore's naturalistic fallacy and open question argument
% Hare's prescriptivism and universalizability
% The logical positivist challenge: are ethical statements meaningful?

\section{Decision Theory and Game Theory}
% Von Neumann--Morgenstern utility
% Arrow's impossibility theorem: the first formal impossibility result for values
% Sen's capability approach: moving beyond scalars

\section{Computational Ethics: Deontic Logic to Machine Ethics}
% Von Wright's deontic logic
% Anderson and Anderson's machine ethics
% The alignment problem: from theoretical to urgent

\section{The Geometric Turn}
% Why differential geometry provides what earlier formalisms lacked
% The analogy with physics: from Newtonian scalars to Einsteinian geometry
% Preview: the moral manifold

\begin{historicalnote}
Arrow's Impossibility Theorem (1951) can be read as an early warning that scalar aggregation of multi-dimensional preferences is fundamentally problematic. The Information Monotonicity Theorem (Chapter~8) generalizes this insight from social choice to moral evaluation.
\end{historicalnote}


% ============================================================
% PART II: MATHEMATICAL SUBSTRATE
% ============================================================
\part{Mathematical Substrate}
\label{part:substrate}

% Chapter 4 ------------------------------------------------
\chapter{Manifolds and Smooth Structure}
\label{ch:manifolds}

\begin{framed}
\noindent\textbf{Learning Objectives.} After completing this chapter, you will be able to:
\begin{enumerate}
    \item Define topological and smooth manifolds from first principles
    \item Construct coordinate charts and atlases
    \item Compute tangent vectors, cotangent vectors, and differentials
    \item Explain why moral evaluation spaces require manifold structure
\end{enumerate}
\end{framed}

\section{Topological Spaces and Continuity}
% Open sets, neighborhoods, continuous maps
% Hausdorff, second-countable --- why these matter for moral spaces

\section{Charts, Atlases, and Smooth Structure}
% Coordinate charts as local parameterizations
% Transition functions and smoothness
% Example: the circle $S^1$ as a manifold

\section{Tangent Spaces and Derivatives}
% Tangent vectors as derivations
% The tangent bundle $T\calM$
% The differential $df$ of a smooth map

\section{Cotangent Spaces and Differential Forms}
% Covectors and the cotangent bundle $T^*\calM$
% 1-forms, 2-forms, the wedge product
% Integration of forms: Stokes' theorem preview

\section{The Moral Manifold: First Construction}
% $\R^9$ as the ambient space for moral evaluation
% The nine dimensions: consequences, rights, fairness, autonomy, privacy,
%   societal impact, virtue, legitimacy, epistemic status
% Why these form a manifold: smoothness of moral tradeoffs within strata

\begin{practitioner}
You do not need to understand every proof in this chapter to use Philosophy Engineering. The key takeaway: a manifold is a space that looks like $\R^n$ locally but can have nontrivial global structure. Moral evaluation lives on such a space because locally, moral tradeoffs are continuous and smooth, but globally, there are boundaries you cannot cross smoothly (rights violations, consent thresholds). Those boundaries are handled in Chapter~6.
\end{practitioner}

\begin{exercise}[A]
Show that the 2-sphere $S^2$ cannot be covered by a single coordinate chart. What does this imply for any attempt to represent all moral evaluations in a single coordinate system?
\end{exercise}

\begin{exercise}[B]
Construct an explicit atlas for the moral manifold restricted to the two-dimensional subspace of (consequences, rights). Identify at least one point where the transition functions are nontrivial.
\end{exercise}

\begin{exercise}[C]
Prove that if $\calM$ is a smooth manifold and $f: \calM \to \R$ is a smooth moral evaluation functional, then $f$ has at least one critical point on any compact submanifold. Interpret this result in terms of moral equilibria.
\end{exercise}


% Chapter 5 ------------------------------------------------
\chapter{Riemannian Geometry and Metrics}
\label{ch:riemannian}

\begin{framed}
\noindent\textbf{Learning Objectives.} After completing this chapter, you will be able to:
\begin{enumerate}
    \item Define Riemannian metrics and compute distances and geodesics
    \item Explain parallel transport and the Levi-Civita connection
    \item Compute curvature tensors and interpret them geometrically
    \item Apply these tools to moral distance and moral geodesics
\end{enumerate}
\end{framed}

\section{Riemannian Metrics}
% Inner products on tangent spaces, the metric tensor $g_{ij}$
% Length of curves, distance function, volume form

\section{Geodesics and the Exponential Map}
% Geodesic equation, shortest paths
% Moral geodesics: the ``least morally costly'' path between states
% The exponential map and normal coordinates

\section{Connections and Parallel Transport}
% The Levi-Civita connection
% Parallel transport around loops: curvature as obstruction
% Physical intuition: transporting a moral standard around a closed path

\section{Curvature}
% Riemann curvature tensor, Ricci tensor, scalar curvature
% Sectional curvature: how moral planes curve
% Positive vs.\ negative curvature: moral spaces that focus vs.\ spread geodesics

\section{The Moral Metric}
% Context-dependent weighting of moral dimensions
% How the metric encodes tradeoff rates between dimensions
% Example: medical triage metric vs.\ content moderation metric

\begin{keyinsight}
The metric tensor $g_{ij}$ on the moral manifold is not arbitrary---it encodes the relative importance and tradeoff rates between moral dimensions in a given context. Changing the context changes the metric, not the manifold. This is why the same moral space can support different but internally consistent ethical frameworks.
\end{keyinsight}


% Chapter 6 ------------------------------------------------
\chapter{Whitney Stratified Spaces}
\label{ch:stratification}

\begin{framed}
\noindent\textbf{Learning Objectives.} After completing this chapter, you will be able to:
\begin{enumerate}
    \item Define stratified spaces and Whitney's conditions (a) and (b)
    \item Explain why smooth manifolds alone are insufficient for moral boundaries
    \item Construct stratified models of consent thresholds and rights violations
    \item State and apply the frontier condition
\end{enumerate}
\end{framed}

\section{Why Smooth Is Not Enough}
% Moral boundaries are sharp: consent/no-consent, legal/illegal
% Smooth manifolds cannot model discontinuities
% The need for controlled singularities

\section{Stratified Spaces}
% Decomposition into smooth strata $\Sigma_0, \Sigma_1, \ldots$
% The frontier condition: $\Sigma_i \cap \overline{\Sigma_j} \neq \emptyset \Rightarrow \Sigma_i \subset \overline{\Sigma_j}$
% Dimension hierarchy: boundaries have lower dimension than regions

\section{Whitney's Conditions}
% Whitney (a): tangent planes of approaching sequences
% Whitney (b): secant lines of approaching sequences
% Why these conditions guarantee ``controlled'' singularities

\section{Moral Stratification}
% Consent thresholds as codimension-1 strata
% Rights violations as lower-dimensional boundaries
% Example: the stratified structure of medical decision-making

\section{Intersection Homology and Moral Topology}
% Goresky--MacPherson intersection homology (conceptual introduction)
% Topological invariants that survive stratification
% What the topology of a moral space tells us about its structure


% Chapter 7 ------------------------------------------------
\chapter{Fiber Bundles, Gauge Theory, and Symmetry}
\label{ch:gauge}

\begin{framed}
\noindent\textbf{Learning Objectives.} After completing this chapter, you will be able to:
\begin{enumerate}
    \item Define fiber bundles, principal bundles, and associated bundles
    \item Explain gauge transformations as changes of local frame
    \item Derive conservation laws from symmetries via Noether's theorem
    \item State the $D_4$ Gauge Group Theorem and interpret it
\end{enumerate}
\end{framed}

\section{Fiber Bundles}
% Total space, base space, fiber, projection
% Trivial vs.\ nontrivial bundles
% The M\"obius band as a nontrivial line bundle

\section{Principal Bundles and Connections}
% Structure group, gauge transformations
% Connections as rules for parallel transport
% Curvature = field strength = failure of parallel transport to close

\section{The Bond--Lens Structure}
% Bonds $\calB$: objective relational data (base space)
% Lenses $\phi$: evaluative perspectives (fiber)
% The principal bundle $\pi: P \to \calM$ in moral context

\section{Group Theory and the $D_4$ Gauge Group}
% The dihedral group $D_4$: symmetries of the square
% Why $D_4$: four fundamental moral symmetries and their compositions
% The $D_4$ Gauge Group Theorem: statement and proof sketch

\section{Noether's Theorem for Ethics}
% Classical Noether: continuous symmetry $\Rightarrow$ conservation law
% The Moral Noether Theorem: re-description invariance $\Rightarrow$ harm conservation
% Conserved normative currents $J^\mu$
% Why harm cannot be made to disappear through relabeling

\section{Gauge-Theoretic Interpretation of Moral Disagreement}
% Different gauges = different ethical frameworks viewing the same moral reality
% Gauge transformations = translations between frameworks
% Gauge-invariant quantities = genuinely objective moral facts

\begin{pitfall}
The analogy between moral gauge theory and physical gauge theory is structural, not metaphysical. We are not claiming that moral phenomena obey the equations of electrodynamics or Yang--Mills theory. The claim is that the \textit{mathematical structure}---bundles, connections, curvature, conservation---applies to both domains because both involve invariance under local re-descriptions.
\end{pitfall}


% ============================================================
% PART III: THE TWO PRINCIPLES
% ============================================================
\part{The Two Principles}
\label{part:principles}

% Chapter 8 ------------------------------------------------
\chapter{The Epistemic Invariance Principle}
\label{ch:eip}

\begin{framed}
\noindent\textbf{Learning Objectives.} After completing this chapter, you will be able to:
\begin{enumerate}
    \item State the EIP formally and explain each component
    \item Prove the non-degeneracy condition
    \item Prove the uncertainty stability condition
    \item Apply EIP to concrete AI system specifications
\end{enumerate}
\end{framed}

\section{Informal Statement}
% ``Differences that do not matter should not make a difference''
% Meaning-preserving redescriptions must not change judgment
% This is not optional: it is a precondition for epistemic competence

\section{Formal Model}
% Input space $X$, output space $Y$, judgment function $f: X \to Y$
% Transformation set $\calT \subseteq \mathrm{Aut}(X)$ of meaning-preserving transforms
% Equivalence relation $\equivX$ generated by $\calT$

\section{The Principle}
\begin{principle}[Epistemic Invariance Principle]
A judgment function $f: X \to Y$ is epistemically invariant with respect to $\calT$ if and only if:
\[
\forall x \in X,\ \forall \tau \in \calT: \quad f(\tau(x)) = f(x)
\]
subject to the non-degeneracy and uncertainty stability conditions.
\end{principle}

\section{Non-Degeneracy}
% A trivial function $f(x) = c$ satisfies invariance vacuously
% Non-degeneracy: $f$ must distinguish at least some genuinely different cases
% Formal condition and proof that it excludes trivial compliance

\section{Uncertainty Stability}
% When $f$ produces uncertainty estimates, invariance extends to them
% $P(f(x) = y \mid \tau(x)) = P(f(x) = y \mid x)$ for all $\tau \in \calT$
% Proof and interpretation

\section{The EIP as a Falsifiable Contract}
% How to test EIP: generate transformation trials, check outputs
% The witness principle: when invariance fails, produce a minimal counterexample
% Connection to software testing and formal verification

\begin{exercise}[A]
A sentiment analysis system assigns scores to product reviews. Define a transformation set $\calT$ appropriate for testing EIP on this system. Include at least five transformations.
\end{exercise}

\begin{exercise}[B]
Prove that if $f$ satisfies EIP with respect to $\calT$ and $g$ satisfies EIP with respect to $\calT$, then $h(x) = (f(x), g(x))$ also satisfies EIP with respect to $\calT$.
\end{exercise}

\begin{exercise}[C]
Construct a function that satisfies EIP with respect to a set $\calT$ but violates non-degeneracy. Then modify it minimally to satisfy both. Prove your construction is correct.
\end{exercise}


% Chapter 9 ------------------------------------------------
\chapter{The Bond Invariance Principle}
\label{ch:bip}

\begin{framed}
\noindent\textbf{Learning Objectives.} After completing this chapter, you will be able to:
\begin{enumerate}
    \item Define bonds and lenses formally
    \item State the BIP and explain its relationship to EIP
    \item Prove the Accountability Theorem
    \item Apply BIP to separate objective structure from subjective evaluation
\end{enumerate}
\end{framed}

\section{Bonds and Lenses}
% Bond $\calB$: the objective relational structure of a moral situation
% Lens $\phi$: an evaluative perspective (consequentialist, deontological, etc.)
% The factorization: moral judgment = $\phi(\calB)$, and $\calB$ must be lens-invariant

\section{Formal Statement}
\begin{principle}[Bond Invariance Principle]
For any bond structure $\calB$ extracted from a moral situation, and any pair of legitimate evaluative lenses $\phi_1, \phi_2$:
\[
\calB(\text{situation under } \phi_1) = \calB(\text{situation under } \phi_2)
\]
The objective relational facts do not change when the evaluative framework changes.
\end{principle}

\section{The Accountability Theorem}
% Statement: If a system satisfies BIP, then for any evaluative disagreement,
% the disagreement can be localized to the lens, not the bond
% Proof
% Consequence: meaningful moral disagreement is always about values, never about facts

\section{BIP and the Bundle Structure}
% Bonds = base space, lenses = fiber
% BIP = well-definedness of the bundle projection
% Gauge transformations change the lens; the base is invariant

\section{Relationship Between EIP and BIP}
% EIP: invariance under re-description (same lens, different surface form)
% BIP: invariance of structure under lens change (different lens, same situation)
% Together: complete invariance --- neither surface form nor evaluative perspective
%   can change the objective relational facts


% Chapter 10 ------------------------------------------------
\chapter{Canonicalization and Enforcement}
\label{ch:canonicalization}

\begin{framed}
\noindent\textbf{Learning Objectives.} After completing this chapter, you will be able to:
\begin{enumerate}
    \item Design canonicalization algorithms for specific domains
    \item Prove the soundness of canonicalization as an enforcement mechanism
    \item Implement approximate compliance with the Bond Index
    \item Build audit artifacts for regulatory and engineering review
\end{enumerate}
\end{framed}

\section{The Mandatory Canonicalization Theorem}
% Post-hoc compliance checking is insufficient for norm-compliant behavior
% Proof: the checking function itself may be representation-dependent
% Consequence: invariance must be enforced *before* evaluation, not verified after

\section{Canonicalization Algorithms}
% Map each equivalence class to a canonical representative
% Requirements: idempotent, deterministic, semantics-preserving
% Domain-specific examples: text, images, structured data

\section{The Bond Index}
% $\beta(f, \calT)$: a quantitative measure of representation independence
% Definition, computation, interpretation
% Threshold values for deployment decisions

\section{Approximate Compliance}
% When perfect invariance is computationally intractable
% $\varepsilon$-compliance: $\|f(\tau(x)) - f(x)\| < \varepsilon$ for all $\tau \in \calT$
% Error propagation through composed transformations

\section{Audit Artifacts}
% What to record: input, transformations applied, outputs, Bond Index, witnesses
% BIP extension fields: bond structure, lens applied, separation verification
% Machine-checkable certificates for regulatory compliance

\section{The EIP Enforcement Pipeline}
% Architecture: Input $\to$ Canonicalize $\to$ Evaluate $\to$ Audit
% Runtime modes: strict rejection, soft flagging, audit-only
% Integration with existing ML pipelines

\begin{casestudy}
\textbf{Medical Triage System.} A hospital deploys an AI system to prioritize emergency department patients. We apply the full EIP enforcement pipeline: define the transformation set (patient name permutation, synonym substitution in symptoms, unit conversion), implement canonicalization (standardize to SNOMED-CT codes), compute the Bond Index, and generate audit artifacts. We show that the pre-canonicalization system has $\beta = 0.73$ and the post-canonicalization system achieves $\beta = 0.98$.
\end{casestudy}


% ============================================================
% PART IV: THE GEOMETRIC ETHICS FRAMEWORK
% ============================================================
\part{Geometric Ethics}
\label{part:geometric-ethics}

% Chapter 11 ------------------------------------------------
\chapter{The Moral Manifold}
\label{ch:moral-manifold}

\begin{framed}
\noindent\textbf{Learning Objectives.} After completing this chapter, you will be able to:
\begin{enumerate}
    \item State and prove the Structured Preservation Theorem
    \item Construct the nine-dimensional moral evaluation space
    \item Define the moral metric for specific application domains
    \item Compute moral distances and geodesics
\end{enumerate}
\end{framed}

\section{The Structured Preservation Theorem}
% Any non-trivial, continuous moral evaluation must be defined on a manifold
% Proof: from the topology of level sets and implicit function theorem
% Corollary: scalar moral scores are necessarily degenerate

\section{The Nine Moral Dimensions}
% Consequences, Rights, Fairness, Autonomy, Privacy,
% Societal Impact, Virtue, Legitimacy, Epistemic Status
% Empirical basis: cross-cultural, cross-temporal validation
% Why nine and not eight or ten: dimensional analysis

\section{Constructing the Moral Metric}
% Context-dependent inner products on $\R^9$
% Example metrics: medical, legal, content moderation, autonomous driving
% Signature of the metric: when dimensions are ``timelike'' vs.\ ``spacelike''

\section{Moral Geodesics and Pathfinding}
% Moral reasoning as $A^*$ search on the moral manifold
% Moral obligations as pre-compiled heuristic functions $h(n)$
% Computational complexity: when pathfinding is tractable

\section{The Information Monotonicity Theorem}
% Formal statement and proof
% Scalar contraction of vector moral assessments is irreversibly lossy
% Quantification: exactly how much information is lost
% Implications for RLHF, reward modeling, and constitutional AI


% Chapter 12 ------------------------------------------------
\chapter{Conservation Laws and the Moral Noether Theorem}
\label{ch:noether}

\begin{framed}
\noindent\textbf{Learning Objectives.} After completing this chapter, you will be able to:
\begin{enumerate}
    \item State and prove the Moral Noether Theorem
    \item Derive conserved normative currents from symmetries
    \item Apply harm conservation to detect moral laundering
    \item Compute the stress-energy tensor for normative systems
\end{enumerate}
\end{framed}

\section{Symmetries and Conservation: The Physical Paradigm}
% Noether's theorem in physics: time invariance $\Rightarrow$ energy conservation
% Translational invariance $\Rightarrow$ momentum conservation
% The deep principle: symmetry is the mother of conservation

\section{The Moral Noether Theorem}
% Re-description invariance of the normative Lagrangian
% Derivation of the conserved normative current $J^\mu$
% Harm conservation: harm cannot be created or destroyed through relabeling

\section{Applications of Harm Conservation}
% Detecting moral laundering: when institutions redistribute harm through renaming
% Corporate restructuring that ``eliminates'' liability without reducing harm
% Algorithmic systems that pass blame through layered abstractions

\section{The Normative Stress-Energy Tensor}
% $T^{\mu\nu}$: how moral ``energy'' and ``momentum'' flow through the normative field
% Physical analogy and formal construction
% Interpretation: pressure, flux, and density in moral systems


% Chapter 13 ------------------------------------------------
\chapter{The $D_4$ Gauge Group and Deontic Structure}
\label{ch:d4}

% The four generators: permission/prohibition, obligation/liberty
% Their compositions and the full $D_4$ structure
% Gauge-invariant observables: what is objectively measurable
% Spontaneous symmetry breaking: how context selects a moral framework


% Chapter 14 ------------------------------------------------
\chapter{The No Escape Theorem}
\label{ch:no-escape}

\begin{framed}
\noindent\textbf{Learning Objectives.} After completing this chapter, you will be able to:
\begin{enumerate}
    \item State the No Escape Theorem precisely
    \item Explain each of the three escape strategies and why they fail
    \item Apply the theorem to evaluate AI safety proposals
    \item Discuss the philosophical implications for moral realism
\end{enumerate}
\end{framed}

\section{Statement of the Theorem}
% Any sufficiently capable AI system is necessarily subject to geometric
% moral constraints that cannot be circumvented through representational manipulation

\section{The Three Escape Strategies}
% Strategy 1: Collapse to scalar --- fails by Information Monotonicity
% Strategy 2: Construct alternative geometry --- fails by uniqueness of Levi-Civita
% Strategy 3: Reject continuity --- fails by computational necessity

\section{Proof}
% Full proof with all lemmas
% Discussion of where each assumption is used

\section{Implications}
% For AI alignment: geometric constraints are not optional design choices
% For metaethics: moral realism without metaphysical commitment
% For regulation: mathematical basis for universal standards


% ============================================================
% PART V: ENGINEERING AND PRACTICE
% ============================================================
\part{Engineering and Practice}
\label{part:engineering}

% Chapter 15 ------------------------------------------------
\chapter{The DEME Architecture}
\label{ch:deme}

\begin{framed}
\noindent\textbf{Learning Objectives.} After completing this chapter, you will be able to:
\begin{enumerate}
    \item Describe the DEME architecture and its components
    \item Explain the evolution from DEME 2.0 (vector) to DEME 3.0 (tensorial)
    \item Implement basic moral evaluation using the ErisML API
    \item Configure DEME for specific application domains
\end{enumerate}
\end{framed}

\section{Architecture Overview}
% DEME = Distributed Ethical Monitoring Engine
% Components: moral evaluator, canonicalizer, auditor, governance layer
% Deployment models: embedded, edge, cloud

\section{From Vectors to Tensors: DEME 2.0 to 3.0}
% DEME 2.0: moral evaluation as vectors in $\R^9$
% Limitations: vectors cannot capture interaction effects between dimensions
% DEME 3.0: full tensorial ethics with metric, connection, curvature
% What tensorial evaluation buys: gauge covariance, harm conservation, geodesic reasoning

\section{ErisML: The Reference Implementation}
% Library architecture and API design
% Core classes: MoralManifold, MetricTensor, GaugeConnection, Evaluator
% Python API walkthrough

\section{Hands-On: Your First Moral Evaluation}
% Step-by-step tutorial with code
% Example: evaluating a content moderation decision
% Computing the Bond Index
% Generating audit artifacts

\begin{lstlisting}[language=Python, caption={Basic moral evaluation with ErisML v3.0}]
import erisml

# Define the moral manifold for content moderation
manifold = erisml.MoralManifold(
    domain="content_moderation",
    dimensions=9
)

# Set context-specific metric
metric = erisml.MetricTensor.from_domain("content_moderation")
manifold.set_metric(metric)

# Evaluate a decision
evaluation = manifold.evaluate(
    action="flag_content",
    context={"content": "...", "user_history": "..."}
)

# Compute Bond Index
beta = erisml.bond_index(
    evaluator=manifold,
    transformations=erisml.standard_transforms("text"),
    sample=test_cases
)
print(f"Bond Index: {beta:.4f}")
\end{lstlisting}


% Chapter 16 ------------------------------------------------
\chapter{Formal Verification of Moral Properties}
\label{ch:verification}

% Type-theoretic encoding of moral invariants
% Model checking for EIP compliance
% Theorem proving for BIP satisfaction
% Certification pipelines for safety-critical deployment


% Chapter 17 ------------------------------------------------
\chapter{Quantum Normative Dynamics}
\label{ch:qnd}

% The quantum formalism applied to normative uncertainty
% Moral superposition: when situations are genuinely ambiguous
% Entanglement: when moral assessments of separate situations are correlated
% Measurement: how observation collapses normative superposition
% Note: most controversial chapter --- clearly labeled as speculative extension


% Chapter 18 ------------------------------------------------
\chapter{Advanced Topics in Moral Geometry}
\label{ch:advanced}

% Stratified gauge theory on singular spaces
% Electrodynamics of value: field-theoretic structure
% Collective agency and multi-agent moral fields
% Moral contraction and the approach to equilibrium
% Open problems and conjectures


% Chapter 19 ------------------------------------------------
\chapter{Case Studies}
\label{ch:case-studies}

\section{Autonomous Vehicle Decision-Making}
% The trolley problem is the wrong framing
% Correct framing: continuous moral evaluation on a stratified manifold
% Implementation with ErisML and ROS/Gazebo
% Results: comparison with scalar reward approaches

\section{Medical Triage and Resource Allocation}
% Multi-dimensional moral evaluation under time pressure
% The triage metric: how to weight dimensions in emergency contexts
% EIP enforcement: ensuring triage is representation-independent
% Audit trail for legal and ethical review

\section{Content Moderation at Scale}
% Cultural variation as gauge transformation
% Invariant core: what is objectively harmful across cultures
% Context-dependent metric: what varies legitimately
% Bond Index benchmarks across languages

\section{Criminal Justice Risk Assessment}
% Representation dependence as a civil rights issue
% EIP as a formal fairness requirement
% Stratification: bright-line rules as moral boundaries
% Audit artifacts for judicial review

\section{Military Autonomous Systems}
% International humanitarian law as geometric constraints
% Distinction, proportionality, precaution as gauge-invariant observables
% The No Escape Theorem applied: why lethal autonomy requires geometric ethics
% DEME deployment architecture for defense applications


% Chapter 20 ------------------------------------------------
\chapter{The Profession of Philosophy Engineering}
\label{ch:profession}

\section{Competency Framework}
% Knowledge, skills, and abilities for PE practitioners
% Three levels: associate, professional, senior
% Mapping to the 16 knowledge areas

\section{Ethical Responsibilities of the Philosophy Engineer}
% The PE practitioner as moral architect
% Responsibilities: transparency, auditability, harm prevention
% When to refuse: ethical red lines for practitioners

\section{Regulatory Landscape}
% EU AI Act and geometric ethics compliance
% IEEE P7000 series and PE alignment
% Future regulatory directions

\section{Open Problems and Research Directions}
% Computational complexity of moral geodesics
% Learning the moral metric from data
% Extending to multi-agent systems
% The measurement problem in quantum normative dynamics
% Connections to category theory and homotopy type theory

\section{An Invitation}
% Philosophy Engineering is not finished --- it is beginning
% How to contribute: open-source, open research, open debate
% The stakes: getting AI alignment right is not optional


% ============================================================
% BACK MATTER
% ============================================================
\backmatter

% --- Appendices ---
\appendix

\chapter{Review of Linear Algebra}
% Vectors, matrices, eigenvalues, inner products, tensor products
% Sufficient for students who need a refresher

\chapter{Review of Point-Set Topology}
% Open/closed sets, compactness, connectedness, quotient spaces
% Sufficient for the manifold chapters

\chapter{Review of Group Theory}
% Groups, subgroups, homomorphisms, representations
% Focus on finite groups and Lie groups used in the text

\chapter{Review of Formal Logic}
% Propositional and predicate logic, proof techniques
% Deontic logic fundamentals

\chapter{ErisML API Reference}
% Complete API documentation for ErisML v3.0
% Class hierarchy, method signatures, parameters
% Configuration options for all supported domains

\chapter{Solutions to Selected Exercises}
% Solutions for all A-level and selected B-level exercises
% Hints for C-level exercises

% --- Bibliography ---
\chapter{Bibliography}

\begin{thebibliography}{99}

\bibitem{bond2026ge}
A.~H.~Bond, \emph{Geometric Ethics: The Mathematical Structure of Moral Reasoning}, v1.14, Department of Computer Engineering, San Jos\'{e} State University, Spring 2026.

\bibitem{bond2026erisml}
A.~H.~Bond, ``ErisML/DEME v3.0: A Library for Governed, Foundation-Model-Enabled Agents with Tensorial Ethics,'' \url{https://github.com/ahb-sjsu/erisml-lib}, 2026.

% --- Foundational Mathematics ---
\bibitem{lee2013}
J.~M.~Lee, \emph{Introduction to Smooth Manifolds}, 2nd ed. Springer, 2013.

\bibitem{docarmo1992}
M.~P.~do Carmo, \emph{Riemannian Geometry}. Birkh\"{a}user, 1992.

\bibitem{nakahara2003}
M.~Nakahara, \emph{Geometry, Topology and Physics}, 2nd ed. CRC Press, 2003.

\bibitem{goresky1988}
M.~Goresky and R.~MacPherson, \emph{Stratified Morse Theory}. Springer, 1988.

\bibitem{noether1918}
E.~Noether, ``Invariante Variationsprobleme,'' \emph{Nachrichten von der Gesellschaft der Wissenschaften zu G\"{o}ttingen}, pp.~235--257, 1918.

% --- Philosophy ---
\bibitem{kant1785}
I.~Kant, \emph{Groundwork of the Metaphysics of Morals}, 1785.

\bibitem{rawls1971}
J.~Rawls, \emph{A Theory of Justice}. Harvard University Press, 1971.

\bibitem{sen1999}
A.~Sen, \emph{Development as Freedom}. Oxford University Press, 1999.

\bibitem{arrow1951}
K.~J.~Arrow, \emph{Social Choice and Individual Values}. Wiley, 1951.

\bibitem{peirce1878}
C.~S.~Peirce, ``How to Make Our Ideas Clear,'' \emph{Popular Science Monthly}, vol.~12, pp.~286--302, 1878.

% --- AI and Alignment ---
\bibitem{christiano2017}
P.~Christiano et al., ``Deep reinforcement learning from human feedback,'' \emph{Advances in Neural Information Processing Systems}, 2017.

\bibitem{bourque2014swebok}
P.~Bourque and R.~E.~Fairley, Eds., \emph{SWEBOK: Guide to the Software Engineering Body of Knowledge}, v3.0. IEEE Computer Society, 2014.

\bibitem{russell2019}
S.~Russell, \emph{Human Compatible: Artificial Intelligence and the Problem of Control}. Viking, 2019.

% Additional references to be added during writing

\end{thebibliography}

% --- Index ---
\printindex

\end{document}
